\section{Related Works}
\label{sec:relatedworks}

\subsection{Movement primitives}
There has been extensive research on movement primitives, including dynamical systems-based approaches such as dynamic movement primitives (DMP)~\cite{saveriano2021dynamic, ijspeert2013dynamical, ijspeert2001trajectory, ijspeert2002learning2, schaal2007dynamics} and stable dynamical systems~\cite{khansari2011learning, neumann2013neural, khansari2014learning}, often formulated to guarantee the stability of the resulting closed-loop system. Additionally, methods that capture diverse motions using probabilistic formulations or manifold-based models have been explored~\cite{calinon2016tutorial, paraschos2013probabilistic, daab2024incremental, noseworthy2020task, lee2023equivariant, lee2024mmp++, beik2021learning}. 

Of particular relevance to our work are \emph{manifold-based models} and \emph{task-conditioned movement primitives}. Recent studies have explored autoencoder-based manifold learning algorithms for learning movement primitives~\cite{noseworthy2020task, beik2021learning, lee2023equivariant, lee2024mmp++}. \cite{beik2021learning} focuses on learning a manifold of {\it configurations}, followed by solving geodesic equations for motion planning. In contrast, \cite{noseworthy2020task, lee2023equivariant, lee2024mmp++} learn a manifold of {\it trajectories}, enabling direct motion sampling. Our approach aligns with the latter, as it also focuses on learning a manifold of trajectories.

Task-Parametrized Gaussian Mixture Model (TP-GMM) can produce new movements for unseen task parameters; however, it primarily addresses parameters defined as coordinate frames, making it less suitable for more general types of inputs such as vision or language~\cite{calinon2016tutorial}.  
Among the manifold-based models mentioned above, \cite{noseworthy2020task, lee2023equivariant} -- except \cite{lee2024mmp++} which does not address task-conditioned motion generation -- adopt conditional autoencoder architectures capable of handling a broader range of task parameters while encoding and generating diverse trajectories, closely aligning with our research goals. However, these methods fall short of capturing the {\it complex} task-motion dependencies within the motion distribution, a limitation we specifically address in this paper.


\subsection{Diffusion and flow matching for imitation learning}
Diffusion and flow matching models -- which fall into the category of probability flow-based models -- have recently gained significant attention as powerful deep generative models that enable efficient and stable training, demonstrating impressive performance on various conditional generation tasks~\cite{song2020denoising, song2020score, ho2022classifier, lipman2022flow, chen2023riemannian, tong2023conditional, lipman2024flow}. 
In robotics, particularly in the context of imitation learning, the use of diffusion and flow matching models for training stochastic {\it policies}, i.e., observation-conditioned action generation models, has recently gained considerable attention. Notable examples include diffusion policy~\cite{chi2023diffusion}, 3D diffusion policy~\cite{ze20243d}, and flow matching policy~\cite{chisari2024learning}, each handling different types of observations such as 2D images and 3D point clouds.

The key difference between recent diffusion or flow matching policies and our method lies in the research objective. While these approaches focus on learning a {\it local policy} that generates low-dimensional {\it local} actions -- both temporally and spatially -- often limited to a few steps of desired configurations, our goal is to develop a {\it global motion planner} capable of producing high-dimensional {\it global} trajectories. This increased dimensionality makes a naive application of diffusion or flow matching methods prone to failure, necessitating motion manifold learning as proposed in this paper. 

Moreover, while local policies can generate actions with high-frequency feedback, the produced actions often lack temporal consistency with previous ones. When actions are defined at the trajectory level, the final state of one action may not smoothly connect with the initial state of the next, leading to non-smooth behavior. For tasks where frequent trajectory re-generation is unnecessary, a global planner is generally better suited for producing long-horizon, smooth motions.